\chapter{练习参考思路}\label{app:answers}
这些练习多涉及设计选择，下面给出可检验的示例思路；实际站点仍以其数据契约和访问规则为准。



\section*{第1章}
第1题：标题和作者可用非空字符串，日期用统一格式，详情链接用绝对 URL；是否允许作者缺失须在采集前约定。第2题：把允许域名设为集合，合成地址后检查主机名，不能只检查原始 \texttt{href}。第3题：空链接跳过；去片段后的重复 URL 由 \texttt{seen} 保序去重，输出仍是两条。

\section*{第2章}
第1题：查询参数会进入 URL，服务端解析得页码 2。第2题：\texttt{raise\_for\_status()} 对 404 抛出 HTTP 错误。第3题：正文可能是错误页或字段不全的 JSON，须结合响应类型与字段验证。

\section*{第3章}
第1题：必填标题缺失时跳过并计数，或抛出结构变化异常；策略须固定。第2题：原选择器将取不到记录，样本测试应失败。第3题：局部查询保留标题与卡片的父子关系，减少错配。

\section*{第4章}
第1题：缺少当前节点前缀可能让每次循环都查询整页，造成错配。第2题：先取 \texttt{data-id}，再与标题和链接组成记录。第3题：两种样本应分别断言记录数为 1 与 0。

\section*{第5章}
第1题：重复游标触发异常，阻止无限循环。第2题：编号 2 只保留一次。第3题：空页可能表示临时错误，\texttt{next} 为空则是接口声明的结束；都应结合接口契约判断。

\section*{第6章}
第1题：跳过并计数，或抛出明确的结构异常，需事先约定。第2题：服务端可能返回错误页面或重定向结果。第3题：可用来源标识与事件编号组成唯一键。

\section*{第7章}
第1题：修改标题后再次运行，同一 \texttt{(source, article\_id)} 键由 \texttt{ON CONFLICT ... DO UPDATE} 更新，记录数仍为 1；\texttt{IF NOT EXISTS} 只处理表已存在，不迁移旧字段。第2题：抓取时间记录观察行为，发布日期记录原文事实。第3题：需要排错或审计时保存脱敏原样本；长期存储应考虑最小化与授权范围。

\section*{第8章}
第1题：并发上限为 1 时五次模拟等待约为串行总时长。第2题：可让整批失败或逐项收集错误，但不得悄悄丢掉失败任务。第3题：总并发为 2 仍可能对同一站点连续发出请求，故需另设站点节流。

\section*{第9章}
第1题：上限为 1 时请求按顺序等待；结果顺序仍由任务列表决定。第2题：保存任务编号与异常类别，按批次策略决定继续或终止。第3题：单独建会话会重复建立连接池，失去复用收益。

\section*{第10章}
第1题：用 \texttt{locator(...).all()} 或计数遍历，逐张取标题。第2题：页面载入事件不保证异步数据已写入 DOM。第3题：JSON 接口依赖较少；浏览器更接近页面行为，但成本和状态更多。

\section*{第11章}
第1题：输出应增加一条记录。第2题：Spider 内校验便于就地理解页面；Pipeline 内校验便于复用与统一统计。第3题：\texttt{get()} 找不到元素时返回 \texttt{None}，必填字段应明确处理。

\section*{第12章}
第1题：缺少或空白 URL 应触发 \texttt{DropItem}。第2题：后续组件依靠返回的记录继续处理。第3题：Feed Export 适合快速导出，SQLite 可借唯一约束与查询处理长期记录。

\section*{第13章}
第1题：示例规则拒绝私有路径，允许公开路径。第2题：超时与暂时性服务错误可有限重试；权限拒绝与字段契约错误应停下来排查。第3题：先推进检查点可能漏掉尚未成功写入的一页。

\section*{第14章}
第1题：至少记录方法、URL、必要参数、状态、响应格式与字段来源。第2题：断言响应类型包含 \texttt{application/json}，且解析后 \texttt{items} 为列表；登录页应触发格式错误。第3题：先比对实际响应，再核对选择器与字段条件。

\section*{第15章}
第1题：首次请求没有会话 Cookie，后续 Session 请求带上本地服务设置的 Cookie。第2题：可记录认证方式、有效期和失败状态；不得记录口令、令牌和完整 Cookie。第3题：还需确认资源授权、服务条款与数据用途。

\section*{第16章}
第1题：依次属于域名解析、TLS、HTTP 和解析校验层。第2题：代理口令也是凭据。第3题：遵从等待指示、降低速率，持续被限流或权限不明时停止并核对规则。

\section*{第17章}
第1题：还应检查实际读到的字节数、格式标识和预期摘要。第2题：相同内容可共享摘要键，但不同来源和授权仍应保留。第3题：OCR 会把相似字符混淆，必须对关键字段抽样核对。

\section*{第18章}
第1题：同一操作可能触发多个请求；时间线和响应内容共同定位因果。第2题：Cookie、认证头、设备标识和个人数据都应检查。第3题：只在获准的测试环境记录脱敏样本，再对照应用行为与接口响应。

\section*{第19章}
第1题：结果为 14。第2题：变量值未知，不能静态折叠。第3题：检查代表性输入输出和异常路径，必要时比较副作用与性能。

\section*{第20章}
第1题：第二页重复编号 2，四条候选最终仍是三条。逐行核对：编号 1 对应 First 与 \texttt{/item/1}，编号 2 对应 Second 与 \texttt{/item/2}，编号 3 对应 Third 与 \texttt{/item/3}。第2题：空标题、空编号或空原始链接在写入前报错，坏页不推进检查点。第3题：在全新 \texttt{build/job.sqlite3} 上模拟第二页失败后执行 \texttt{--resume}；报表应为恢复一页、抓取一页、共三条唯一记录。

\section*{第21章}
第1题：共享去重约束让重复投递只产生一条结果。第2题：两台机器各有自己的内存集合，无法相互保证唯一。第3题：随机抽样核对标题、链接与来源页的对应率。

\section*{第22章}
第1题：零并发不满足配置下界；\texttt{NaN} 与正无穷超时不满足有限正数条件，均须报错。第2题：先看请求与错误统计，再核对响应样本和解析器。第3题：日志说明程序发生了什么，抽样说明数据是否符合目标；运行脚本的非零退出码用于调度器报警。
